%% Extraído de docs/TCC_documento/Textuais/Revisao/01_Helice.tex (linhas 87-130)
    \begin{tikzpicture}
        \begin{axis}[
            axis lines = left,
            xlabel = {Rotação Mecânica ($\omega$) [pu]},
            ylabel = {Torque ($T$) [pu]},
            xmin=0, xmax=1.5,
            ymin=0, ymax=1.2,
            xtick={0, 0.5, 1.0, 1.5},
            ytick={0, 0.5, 1.0, 1.5},
            grid=both,
            grid style={dotted, gray},
            legend style={at={(0.95,0.95)}, anchor=north east, nodes={scale=0.8}},
            height=7cm,
            width=12cm,
            legend style={
                at={(0.5,0)}, 
                anchor=north, yshift=-1.1cm, legend columns=2,
                nodes={scale=0.8, text=black}, % Texto sempre preto
                cells={anchor=west}
            },
        ]
        
        % Curva do Motor em PU
        % Zona 1: Torque constante até a rotação base (w = 1)
        \addplot[domain=0:1.0, blue, thick, name path=motor2, forget plot] {1};
        % Zona 2: Enfraquecimento de campo (T decai com 1/w)
        \addplot[domain=1.0:1.5, blue, thick, name path=motor] {1/x}; 
        \addlegendentry{Torque Motor}
        % Curva da Hélice em PU
        % Modelo proporcional a w^2. O coeficiente (ex: 0.6) define o "peso" da hélice.
        \addplot[domain=0:1.5, red, thick, name path=helice] {2*x^2};
        \addlegendentry{Torque Hélice}
        % Encontra a intersecção automaticamente
        \path [name intersections={of=motor2 and helice, by={point1}}];
        
        % Marca o ponto de intersecção
        \fill (point1) circle (2.5pt) node[above left=0.1cm, fill=white, inner sep=1pt, draw, rounded corners, scale=0.8] {Ponto de Operação};

        % Linhas tracejadas de auxílio para os eixos (usando projeção ortogonal |- e -| para evitar erros de compilação)
        \draw[dotted, thick] (point1) -- (point1 |- {axis cs:0,0}) node[below, fill=white, inner sep=2pt] {\scriptsize $\omega_{op}$};
        \draw[dotted, thick] (point1) -- (point1 -| {axis cs:0,0}) node[left, fill=white, inner sep=2pt] {\scriptsize $T_{op}$};

        \end{axis}
    \end{tikzpicture}
